
\subsubsection{Data Collection}

Metadata fields were collected from three repositories:

\begin{itemize}
\item \textbf{HuggingFace Hub}: 150 fields (75 from datasets with Open Datasheets structured markup, 75 from datasets predating Open Datasheets adoption)
\item \textbf{OpenML}: 100 fields from datasets with structured XML metadata
\item \textbf{UCI Machine Learning Repository}: 50 fields from plain text README files
\end{itemize}

Total sample size: 300 fields. Fields included field names and example values (e.g., ''license: MIT'', ''description: This dataset contains...'').

\subsubsection{Annotation Protocol}

The original experimental plan specified that three expert annotators would independently label all 300 fields using the two-tier lifecycle taxonomy (General Information versus Responsible AI). However, the actual implementation used content-based heuristics to simulate annotations rather than human expert labeling. This limitation is acknowledged in the validation reports, which note that human annotations would be required for valid inter-annotator agreement measurement.

Despite this limitation, the lifecycle labels assigned were used for subsequent analyses. The class distribution was: 275 General Information fields (91.7\%), 25 Responsible AI fields (8.3\%), yielding an 11:1 ratio.

\subsubsection{Embedding Model}

All metadata fields were embedded using the sentence-transformers model \texttt{all-MiniLM-L6-v2} \citep{Reimers2019}, a 22-million-parameter model producing 384-dimensional embeddings. The model was used in frozen form without fine-tuning to isolate distributional signals present in the pretrained representations.

\subsubsection{Experimental Design}

The study followed a two-stage validation protocol:

\textbf{Stage 1 (H-E1): Signal Existence Validation}
\begin{itemize}
\item Inter-annotator agreement measurement (target: Cohen's $\kappa \geq$ 0.60)
\item Linear probe training on 60 scaffolded HuggingFace samples
\item Cross-repository probe testing (target accuracy $\geq$ 0.75)
\end{itemize}

\textbf{Stage 2 (H-M-integrated): Unsupervised Clustering Test}
\begin{itemize}
\item K-means clustering (k=2) on all 300 embedded fields
\item Baseline comparisons: permutation, LDA topic modeling (2 topics), lexical keyword matching
\item Repository stratification analysis
\item Scaffolding effect analysis (scaffolded versus unscaffolded HuggingFace samples)
\end{itemize}

Success criteria for unsupervised clustering: NMI > 0.60 and baseline improvement $\geq$ 0.15.

\subsubsection{Evaluation Metrics}

\begin{itemize}
\item \textbf{Normalized Mutual Information (NMI)}: Cluster-label agreement, range [0, 1], normalized for size imbalance
\item \textbf{Cohen's $\kappa$}: Inter-annotator agreement
\item \textbf{Probe accuracy}: Linear probe classification accuracy on held-out data
\end{itemize}

